\scanpage{192}
% manuscript position: top
\mspage{108}
% manuscript position: 1
\uline{Dém}. On a les implications $\pm$ tautologiques
% manuscript position: 2–3
\diagram{\begin{tikzcd}[ampersand replacement=\&,column sep=small,row sep=tiny]
\deleted{\unsure{a}}\arrow[r,phantom,"\deleted{\unsure{\Longrightarrow}}" description]\arrow[dd,phantom,"\deleted{\unsure{\Downarrow}}" description] \& a'\arrow[r,Rightarrow] \& a'' \& {}\\
{(\alpha)}\quad a\arrow[ur,Rightarrow]\arrow[dr,Rightarrow] \& {} \& {} \& {}\\
\deleted{b} \& b\arrow[r,Rightarrow] \& c'\arrow[r,Rightarrow] \& c^{\deleted{\unsure{\prime}}}\text{,}
\end{tikzcd}}
% manuscript position: 4–9
où $a'\Rightarrow a''$ provient de l'inclusion canonique $F(x)=\operatorname{Hom}_M(x,F)\mathrel{\lhook\joinrel\relbar}\operatorname{Ob} C_{/F}$, $a\Rightarrow b$ s'obtient en prenant $I=C_{0/F}$ et notant que $I$ est stable par limites inductives \deleted{grandes} \inserted{de} cardinal $\leq\pi_0$, ($C$ l'étant), donc est \unsure{grde|grdes} devant $\pi_0$.
% manuscript position: left/6–10
% slanted
\margin{on n'a pas besoin de cette stabilité, de \unsure{ttes} façons $C_{0/F}$ est \unsure{grde} devant $\pi_0$, \unsure{en vertu de|vu} 4.4.3 p. 80}

% manuscript position: 9–11
On a $b)\Rightarrow c'$ en vertu du fait que la sous-catégorie \inserted{str. pleine} $\unsure{\widehat M_\pi|\widehat C_\pi}$ de $\unsure{\widehat C|\widehat M}$ définie par
% manuscript position: 12–13
\[
\left\{\begin{aligned}
\widehat M_\pi={}&\text{sous-cat. \inserted{str.} pleine de }\widehat C\text{ formée des}\\
&F\in M\ \deleted{\text{\unsure{tels}}}\text{ qui sont }\pi\text{-accessibles}
\end{aligned}\right.\tag*{(\unsure{4.5.2.1|4.2.2.1})}
\]
% manuscript position: 14–19
est stable par $\varinjlim$ de cardinal $\leq\pi$ (SGA 4 I 9.9). Enfin $c'\Rightarrow c^{\unsure{\prime\prime|\prime}}$ provient du fait que le \deleted{\unsure{lim}} foncteur d'inclusion $M\hookrightarrow\widehat C$ est pl. fidèle et \uline{$\pi_0$-accessible}, i.e. commute aux limites inductives grandes devant $\pi_0$, et a fortiori à celles grandes devant $\pi$.

% manuscript position: 20–22
\deleted{Donc pour \unsure{achever}} \deleted{\unsure{la}} Par les indications du diagramme ci-dessus, il suffit donc de prouver qu'on a
% manuscript position: left/22
\margin{$\unsure{2|Z}$}
% manuscript position: 23
\[
c\overset{\inserted{\text{n'est prouvé (et n'est vrai) que si }\pi^{\pi_0}=\pi}}{\Rightarrow} a
\qquad\text{et}\qquad a''\Rightarrow a.
\]
% manuscript position: 24–26
\deleted{On \unsure{va|en} d'abord \unsure{prouver|prouvons} \deleted{\unsure{$a\Rightarrow$}} $a''\Rightarrow a$, ce qui implique déjà \unsure{l'équivalence|l'égalisation} des}
% manuscript position: foot
% archive sheet number: 191
